\chapter{Boundary Conditions and the Helmholtz Decomposition}
\label{ch:bc}

\chapterorient{The assembly machinery of the last three chapters handed us a
linear system $\mat{K}\vx=\vb$ built term by term from the weak form. But two
things stand between that raw system and a physically correct answer. First, the
boundary conditions: the weak form of \cref{ch:weak} split them into
\emph{essential} and \emph{natural}, and the essential half---the voltages on
conductors, the tangential field on metal walls---was deferred to ``the discrete
setting.'' This is that setting. We make essential conditions concrete as an
operation on the assembled matrix: identify the constrained degrees of freedom,
zero their rows and columns, and lift any nonzero boundary data into the
right-hand side. Second, the fields the solver produces may not be physical: the
curl--curl operator has an enormous null space of gradient fields, and an
eigensolver left to its own devices will happily return them. The cure is a
\emph{projection}---the discrete Helmholtz decomposition---that strips the
gradient part off a field and keeps only the divergence-free remainder. Both are
\emph{finishing operations}: post-compositions that act on the assembled system,
cleanly separable from the assembly fold that produced it. They close Part~II.}

\section{Two finishing operations}

The previous chapters built an operator. Out of the mesh, the materials, and a
short list of weak-form terms $(Q,\diffop)$, the assembly fold of
\cref{ch:assembly} produced a sparse matrix $\mat{K}$ and, where there is a
source, a load vector $\vb$. We are nearly ready to solve $\mat{K}\vx=\vb$. Two
operations remain, and this chapter is about both. They are different in
character but share a structural signature, and naming that shared signature is
the first order of business.

The first operation \emph{imposes the essential boundary conditions}. The weak
form of \cref{ch:weak} taught us that boundary conditions come in two flavors.
\emph{Natural} (Neumann/impedance) conditions take care of themselves: they enter
as boundary integrals during assembly, and a homogeneous one contributes nothing
at all---you get it for free by doing nothing. \emph{Essential} (Dirichlet)
conditions are the work. They fix the value of the solution on part of the
boundary---a conductor held at a known voltage, the tangential electric field
forced to zero on a perfectly conducting wall---and that constraint must be
\emph{built into the discrete system by hand}, after assembly, by operating on
the rows and columns of $\mat{K}$ that belong to the constrained unknowns.

The second operation \emph{projects a field onto the physical subspace}. Some of
the operators we assemble---the curl--curl operator above all---are singular:
they annihilate an entire subspace of fields, the gradients. A solver handed such
an operator can drift into that null space and return a field that is
mathematically a solution but physically meaningless. The remedy is to project
the iterate, at each step, onto the subspace orthogonal to the null space---the
\emph{divergence-free} fields. That projection is the discrete realization of a
classical theorem, the Helmholtz decomposition, and it is the second half of the
chapter.

\begin{figure}[t]
\centering
\begin{tikzpicture}[node distance=9mm]
  \node[stage, text width=26mm, align=center] (asm)
    {\textbf{assemble}\\[1pt]\footnotesize fold over terms\\\scriptsize(\cref{ch:assembly})};
  \node[product, right=12mm of asm, text width=22mm, align=center] (raw)
    {raw system\\[1pt]\footnotesize $\mat{K}\vx=\vb$};
  \node[opbox, right=12mm of raw, text width=30mm] (bc)
    {\textbf{impose essential BC}\\[1pt]\scriptsize zero rows/cols, lift RHS};
  \node[opbox, below=11mm of bc, text width=30mm] (proj)
    {\textbf{divergence-free project}\\[1pt]\scriptsize strip the gradient part};
  \node[stage, right=12mm of bc, text width=20mm, align=center, fill=simBgGreen, draw=simGreen] (solve)
    {\textbf{solve}\\[1pt]\footnotesize (\cref{ch:krylov})};
  \draw[flow] (asm) -- (raw);
  \draw[flow] (raw) -- (bc);
  \draw[flow] (bc) -- (solve);
  \draw[flow] (proj.east) -- ++(6mm,0) |- (solve.west);
  \draw[refedge] (raw.south) |- (proj.west);
  \node[font=\scriptsize\itshape, simGray, below=1mm of asm.south] {the fold};
  \node[font=\scriptsize\itshape, simRust] at ($(bc.north)+(0,5mm)$) {finishing operations (this chapter)};
\end{tikzpicture}
\caption[Where the finishing operations sit]{The two finishing operations of this
chapter sit \emph{between} the assembly fold and the solve. Essential-boundary
elimination acts on the assembled operator $\mat{K}$ and its right-hand side
$\vb$; the divergence-free projection acts on field iterates during (or before)
the solve. Both are \emph{post-compositions}: separable from the assembly that
produced the raw system, they neither inspect nor depend on which weak-form terms
were folded in.}
\label{fig:finishing}
\end{figure}

What unifies them is that both are \emph{post-compositions} on the assembled
system (\cref{fig:finishing}). Neither is part of the assembly fold. Essential
elimination consumes the finished operator $\mat{K}$ and does not care which
weak-form terms went into it; the divergence-free projection consumes a field and
a few auxiliary operators and likewise does not crack open the assembly. This
separability is not an accident of implementation---it is the reason the
simulator can assemble once and then apply boundary conditions and projections as
modular, swappable stages. We will make the claim precise for each operation in
turn.

\begin{keyidea}
The assembled system $\mat{K}\vx=\vb$ is not yet ready to solve. Two
\emph{finishing operations} stand between assembly and solution, and both are
\emph{separable post-compositions} on the assembled operator---not steps inside
the assembly fold. \textbf{(1)~Essential-boundary elimination} forces the
solution to take prescribed values on the Dirichlet boundary, by zeroing the
constrained rows and columns of $\mat{K}$ and lifting the boundary data into
$\vb$. \textbf{(2)~Divergence-free projection} removes the unphysical gradient
component of a field, the discrete Helmholtz decomposition, so that singular
operators like curl--curl are solved on the subspace where they are well-posed.
\end{keyidea}

\section{Essential boundary conditions on the discrete system}
\label{sec:essential}

\subsection{The essential degrees of freedom}

Recall from \cref{ch:galerkin} that the discrete unknown $\vx$ is a vector of
\emph{degrees of freedom} (DoFs): one number per basis function, the coefficient
that scales that basis function in the finite-element expansion of the solution.
For an $\Hone$ (nodal) space, a DoF is the value of the potential at a mesh
vertex; for an $\Hcurl$ (Nédélec) space, a DoF is the circulation of the field
along a mesh edge. The full set of DoFs indexes the rows and columns of
$\mat{K}$.

An essential boundary condition pins some of those DoFs to known values. On a
conductor held at voltage $V$, every $\Hone$ DoF whose vertex lies on that
conductor must equal $V$. On a perfectly conducting (PEC) wall, the tangential
electric field vanishes, which means every $\Hcurl$ DoF whose edge lies in that
wall must equal zero (the edge circulation of a field with no tangential
component is zero). In both cases a geometrically identifiable subset of the DoFs
is \emph{constrained}: its values are dictated by the boundary condition, not by
the equation.

\begin{definition}[Essential and free DoFs]
\label{def:essdof}
Let the full DoF index set be $\{0,1,\dots,N-1\}$. The \emph{essential} (or
constrained) DoF set $E$ is the set of indices whose values are prescribed by an
essential boundary condition. The \emph{free} DoF set is its complement
$F = \{0,\dots,N-1\}\setminus E$. The two are disjoint and exhaust the index set:
$E\cup F = \{0,\dots,N-1\}$, $E\cap F=\emptyset$.
\end{definition}

The map from boundary geometry to the essential set $E$---walk the boundary
elements carrying the essential attribute, collect every DoF they touch---is a
purely combinatorial pre-pass, run once before any matrix is touched. We treat
$E$ as given. \Cref{fig:partition} shows the partition on a small mesh: the
boundary DoFs (essential, dark) form a thin shell, and the interior DoFs (free,
light) are everything else.

\begin{figure}[t]
\centering
\begin{tikzpicture}[scale=1.0]
  % domain
  \fill[simBgBlue] (0,0) rectangle (5,3);
  % grid lines
  \foreach \x in {0,1,2,3,4,5}{\draw[simGray!40] (\x,0)--(\x,3);}
  \foreach \y in {0,1,2,3}{\draw[simGray!40] (0,\y)--(5,\y);}
  % boundary (essential) nodes: dark, on the perimeter
  \foreach \x in {0,1,2,3,4,5}{
    \fill[simBlue] (\x,0) circle (2.6pt);
    \fill[simBlue] (\x,3) circle (2.6pt);}
  \foreach \y in {1,2}{
    \fill[simBlue] (0,\y) circle (2.6pt);
    \fill[simBlue] (5,\y) circle (2.6pt);}
  % interior (free) nodes: light
  \foreach \x in {1,2,3,4}{
    \foreach \y in {1,2}{
      \fill[simRust] (\x,\y) circle (2.6pt);}}
  % labels
  \node[simBlue, font=\small, below left] at (0,0) {};
  \node[font=\small\bfseries, simBlue] at (2.5,3.45) {essential DoFs $E$ (boundary)};
  \node[font=\small\bfseries, simRust] at (2.5,1.5) {};
  \node[font=\small\bfseries, simRust, anchor=west] at (5.4,1.5) {free DoFs $F$ (interior)};
  \draw[refedge] (5.35,1.5) -- (4,1.5);
  % a constrained value callout
  \node[font=\scriptsize, simBlue, anchor=east] at (-0.3,1.5) {$x_i=V$};
  \draw[refedge] (-0.3,1.5) -- (0,1);
\end{tikzpicture}
\caption[The free/essential DoF partition]{The degrees of freedom on a small
$\Hone$ mesh split into two disjoint groups. The \emph{essential} DoFs $E$ (blue,
on the boundary) carry prescribed values---here a conductor voltage $x_i=V$---and
are not solved for. The \emph{free} DoFs $F$ (rust, interior) are the genuine
unknowns. Imposing the essential condition means reducing the system to the free
block, with the influence of the prescribed boundary values moved to the
right-hand side.}
\label{fig:partition}
\end{figure}

\subsection{Scalar value versus tangential trace}
\label{sec:scalarvstrace}

The two physical pictures---a conductor at fixed voltage and a PEC wall forcing
tangential $\vE$ to zero---feel different, but they produce essential DoF sets in
exactly the same way, once you remember what a DoF \emph{is} in each space. The
difference is entirely in the trace operator the boundary condition constrains.

In an $\Hone$ (nodal) space the relevant trace is the boundary \emph{value} of the
function. A DoF is a point value at a vertex, so ``$u=u_0$ on the conductor''
constrains every DoF whose vertex lies on the conductor. In an $\Hcurl$ (Nédélec)
space the relevant trace is the \emph{tangential component} of the field on the
boundary surface. A DoF is an edge circulation $\int_e \vu\cdot\mathrm{d}\boldsymbol{\ell}$,
and an edge lying \emph{in} a boundary face sees only the tangential part of the
field; so ``tangential $\vE=0$ on the PEC wall'' constrains every edge DoF whose
edge lies in that wall. In both cases the recipe is the same: walk the boundary
entities carrying the essential attribute, collect the DoFs they own, and add them
to $E$. That the constrained quantity is a point value in one space and a
tangential circulation in the other never surfaces in the elimination---by the
time \texttt{eliminate\_essential\_bc} runs, $E$ is just a list of indices.

\begin{figure}[t]
\centering
\begin{tikzpicture}[scale=1.0]
  % --- H1 nodal: value at vertices ---
  \begin{scope}[shift={(0,0)}]
    \node[font=\small\bfseries, simBlue] at (1.5,2.9) {$\Hone$: nodal value};
    \draw[simGray!50] (0,0) rectangle (3,2.2);
    \draw[simGray!40] (1.5,0)--(1.5,2.2);
    \draw[simGray!40] (0,1.1)--(3,1.1);
    % bottom boundary face highlighted
    \draw[simBlue, line width=2pt] (0,0)--(3,0);
    % boundary vertices (essential): on the bottom edge
    \foreach \x in {0,1.5,3}{\fill[simBlue] (\x,0) circle (3pt);}
    % interior vertices
    \fill[simRust] (1.5,1.1) circle (3pt);
    \foreach \x in {0,1.5,3}{\fill[simRust] (\x,2.2) circle (2.4pt);}
    \foreach \x in {0,3}{\fill[simRust] (\x,1.1) circle (2.4pt);}
    \node[font=\scriptsize, simBlue, below] at (1.5,-0.15) {value $u=u_0$ pinned};
    \node[font=\scriptsize, simRust] at (2.3,1.35) {free node};
  \end{scope}
  % divider
  \draw[simGray,densely dashed] (4.0,-0.5)--(4.0,3.0);
  % --- Hcurl edge: tangential circulation ---
  \begin{scope}[shift={(5.2,0)}]
    \node[font=\small\bfseries, simGreen] at (1.5,2.9) {$\Hcurl$: edge circulation};
    \draw[simGray!50] (0,0) rectangle (3,2.2);
    \draw[simGray!40] (1.5,0)--(1.5,2.2);
    \draw[simGray!40] (0,1.1)--(3,1.1);
    % bottom boundary face highlighted
    \draw[simGreen, line width=2pt] (0,0)--(3,0);
    % boundary edges (essential): the two bottom edges -> arrows along them
    \draw[-{Stealth[length=2mm]}, simGreen, very thick] (0.1,0)--(1.4,0);
    \draw[-{Stealth[length=2mm]}, simGreen, very thick] (1.6,0)--(2.9,0);
    % an interior edge (free)
    \draw[-{Stealth[length=2mm]}, simRust, very thick] (1.5,1.2)--(1.5,2.1);
    \node[font=\scriptsize, simGreen, below] at (1.5,-0.15) {tangential $\vu\cdot\boldsymbol{\ell}=0$};
    \node[font=\scriptsize, simRust, anchor=west] at (1.6,1.6) {free edge};
  \end{scope}
\end{tikzpicture}
\caption[Essential DoFs for the scalar and vector spaces]{The essential DoF set is
collected the same way in both function spaces; only the constrained trace
differs. In an $\Hone$ nodal space (left) a DoF is a point value at a vertex, so a
Dirichlet value condition $u=u_0$ pins every vertex DoF on the boundary face
(blue). In an $\Hcurl$ Nédélec space (right) a DoF is the circulation along an
edge, so a tangential-trace condition (PEC wall, tangential $\vE=0$) pins every
\emph{edge} DoF lying in the boundary face. In both cases the constrained DoFs are
collected into $E$ and eliminated identically; the elimination never knows which
kind of trace was constrained.}
\label{fig:trace}
\end{figure}

\Cref{fig:trace} sets the two side by side. The practical consequence is uniformity:
the same \texttt{eliminate\_essential\_bc} and \texttt{eliminate\_rhs} serve the
electrostatic (nodal) and magnetostatic/eigenmode (edge) pipelines without
modification, because the only space-specific work---deciding which DoFs are
constrained---happens in the combinatorial pre-pass that builds $E$, upstream of
the matrix surgery.

\subsection{Why this is a post-composition, not part of assembly}

Here is the central structural point, and it echoes the slogan of
\cref{fig:finishing}. The assembly fold of \cref{ch:assembly} knows nothing about
boundary conditions. It walks the list of weak-form terms, and for each term it
adds an element contribution into the global matrix; it would produce the very
same $\mat{K}$ whether or not any DoF is later constrained. Essential elimination
is applied \emph{afterward}, to the finished matrix, as a separate operation:
\[
  \mat{K}_{\mathrm{BC}} \;=\; \texttt{eliminate\_essential\_bc}\bigl(\mat{K},\,E,\,\text{policy}\bigr),
  \qquad
  \mat{K} = \texttt{assemble}(\text{space},\,\text{terms}).
\]
The elimination consumes $\mat{K}$ as an opaque square operator together with the
index set $E$. It never asks which terms were folded to build $\mat{K}$, nor in
what representation it is stored. This is what it means to be a \emph{separable
post-composition}: it factors cleanly behind the assembly, and the two compose
without either needing to know the other's internals.

\begin{pitfall}
A natural instinct is to ``bake the boundary conditions into assembly''---to skip
the constrained DoFs while folding, so that $\mat{K}$ comes out already reduced.
Resist it. Keeping elimination as a separate post-composition is what lets the
\emph{same} assembled operator feed a Dirichlet solve, a Neumann solve, or an
eigenvalue problem, each with its own essential set and its own diagonal policy.
It also lets us reason about the two stages independently: assembly is a fold;
elimination is a projection. Fusing them couples two concerns that have no
business being coupled.
\end{pitfall}

\subsection{Eliminating the essential DoFs}

Now the operation itself. Partition the unknown vector and the matrix along the
free/essential split of \cref{def:essdof}. Writing the unknown as
$\vx=(\vx_F,\vx_E)$ and ordering rows and columns so that free DoFs come first,
the assembled system $\mat{K}\vx=\vb$ has the block structure
\begin{equation}
  \begin{bmatrix} \mat{K}_{FF} & \mat{K}_{FE} \\[2pt]
                  \mat{K}_{EF} & \mat{K}_{EE} \end{bmatrix}
  \begin{bmatrix} \vx_F \\[2pt] \vx_E \end{bmatrix}
  =
  \begin{bmatrix} \vb_F \\[2pt] \vb_E \end{bmatrix}.
  \label{eq:blocksys}
\end{equation}
The block $\mat{K}_{FF}$ couples free DoFs to free DoFs---the genuine interior
physics. The off-diagonal blocks $\mat{K}_{FE}$ and $\mat{K}_{EF}$ couple free to
essential. The essential values $\vx_E$ are \emph{known}; only $\vx_F$ is to be
solved for.

The first row block of \cref{eq:blocksys} reads
$\mat{K}_{FF}\vx_F + \mat{K}_{FE}\vx_E = \vb_F$. Since $\vx_E$ is known, move its
contribution to the right:
\begin{equation}
  \mat{K}_{FF}\,\vx_F \;=\; \vb_F - \mat{K}_{FE}\,\vx_E.
  \label{eq:reduced}
\end{equation}
This reduced system is the honest problem: a smaller, square, well-posed system
on the free DoFs alone, with the known boundary data appearing as a forcing term
$-\mat{K}_{FE}\vx_E$ on the right. If we could re-index the unknowns at will we
would simply solve \cref{eq:reduced} and be done.

In practice the simulator keeps the full $N\times N$ index space---re-indexing a
sparse matrix is expensive and entangles boundary handling with storage---and
realizes the same reduction \emph{in place}. The operation
\texttt{eliminate\_essential\_bc} produces a fresh operator equal to $\mat{K}$
with the essential rows and columns zeroed and the eliminated diagonal set by a
\emph{policy}:
\begin{equation}
  \texttt{eliminate\_essential\_bc}(\mat{K},E,\text{policy})
  \;=\;
  \begin{bmatrix} \mat{K}_{FF} & \mathbf{0} \\[2pt]
                  \mathbf{0} & \mat{D} \end{bmatrix},
  \qquad
  \mat{D} =
  \begin{cases}
    \mat{I}_E & \text{(policy } \mathtt{DIAG\_ONE}),\\
    \mathbf{0}_E & \text{(policy } \mathtt{DIAG\_ZERO}).
  \end{cases}
  \label{eq:elim}
\end{equation}
Three things happen. The off-diagonal coupling blocks $\mat{K}_{FE}$ and
$\mat{K}_{EF}$ are zeroed---this decouples the essential DoFs from the free ones,
so the essential equations no longer pollute the free block and vice versa. The
free--free block $\mat{K}_{FF}$ is preserved exactly---the interior physics is
untouched. And the essential diagonal is overwritten with the identity or with
zero, depending on the policy. \Cref{fig:elim} shows the matrix before and after.

\begin{figure}[t]
\centering
\begin{tikzpicture}[scale=0.78]
  % ---- before ----
  \begin{scope}[shift={(0,0)}]
    \node[font=\small\bfseries] at (1.9,4.3) {before: $\mat{K}$};
    % FF block (large, top-left)
    \fill[simBgBlue] (0,1.2) rectangle (2.8,4);
    \node[simBlue,font=\small] at (1.4,2.6) {$\mat{K}_{FF}$};
    % FE block (top-right strip)
    \fill[simBgRust] (2.8,1.2) rectangle (3.8,4);
    \node[simRust,font=\scriptsize] at (3.3,2.6) {$\mat{K}_{FE}$};
    % EF block (bottom-left strip)
    \fill[simBgRust] (0,0) rectangle (2.8,1.2);
    \node[simRust,font=\scriptsize] at (1.4,0.6) {$\mat{K}_{EF}$};
    % EE block (bottom-right)
    \fill[simBgGray] (2.8,0) rectangle (3.8,1.2);
    \node[simGray,font=\scriptsize] at (3.3,0.6) {$\mat{K}_{EE}$};
    \draw[simInk,thick] (0,0) rectangle (3.8,4);
    \draw[simGray] (2.8,0)--(2.8,4);
    \draw[simGray] (0,1.2)--(3.8,1.2);
  \end{scope}
  % arrow
  \node[font=\Large] at (5,2) {$\longrightarrow$};
  \node[font=\scriptsize,simGray,above] at (5,2.2) {eliminate};
  % ---- after ----
  \begin{scope}[shift={(6.4,0)}]
    \node[font=\small\bfseries] at (1.9,4.3) {after: $\mat{K}_{\mathrm{BC}}$};
    \fill[simBgBlue] (0,1.2) rectangle (2.8,4);
    \node[simBlue,font=\small] at (1.4,2.6) {$\mat{K}_{FF}$};
    % zeroed FE
    \fill[white] (2.8,1.2) rectangle (3.8,4);
    \node[simGray,font=\scriptsize] at (3.3,2.6) {$\mathbf{0}$};
    % zeroed EF
    \fill[white] (0,0) rectangle (2.8,1.2);
    \node[simGray,font=\scriptsize] at (1.4,0.6) {$\mathbf{0}$};
    % diagonal D
    \fill[simBgGreen] (2.8,0) rectangle (3.8,1.2);
    \node[simGreen,font=\scriptsize] at (3.3,0.6) {$\mat{D}$};
    \draw[simInk,thick] (0,0) rectangle (3.8,4);
    \draw[simGray] (2.8,0)--(2.8,4);
    \draw[simGray] (0,1.2)--(3.8,1.2);
    % diagonal hint inside D
    \draw[simGreen,thick] (2.85,1.15)--(3.75,0.05);
  \end{scope}
\end{tikzpicture}
\caption[Essential-BC elimination on the matrix blocks]{Essential elimination as a
block operation. The assembled $\mat{K}$ (left) couples free and essential DoFs
through the off-diagonal blocks $\mat{K}_{FE}$, $\mat{K}_{EF}$. Elimination
(right) \emph{zeroes} those coupling blocks and the essential--essential block,
\emph{preserves} the free--free block $\mat{K}_{FF}$ exactly, and writes a chosen
diagonal $\mat{D}$ on the essential block---the identity $\mat{I}_E$ under the
\texttt{DIAG\_ONE} policy (so each essential equation reads $x_i=b_i$) or zero
under \texttt{DIAG\_ZERO}. The interior physics in $\mat{K}_{FF}$ is never
touched.}
\label{fig:elim}
\end{figure}

The two policies serve two purposes. \texttt{DIAG\_ONE} sets the essential
diagonal to the identity, so each essential row becomes the trivial equation
$1\cdot x_i = b_i$: the matrix is non-singular on the essential block, and the
prescribed value rides in through the right-hand side. This is the
\emph{solve-side} convention---it is what you use when the eliminated operator
will be inverted by a linear solver. \texttt{DIAG\_ZERO} leaves the essential
diagonal at zero, the \emph{energy/mass-block} convention used when the essential
block must contribute no spurious unit eigenvalues, as when assembling the blocks
of a generalized eigenvalue problem.

There is a clean algebraic way to see what elimination is. Let $\mat{P}_F$ be the
diagonal $0/1$ matrix that projects onto the free DoFs ($1$ on free indices, $0$
on essential). Then the \texttt{DIAG\_ZERO} elimination is exactly the two-sided
projection
\begin{equation}
  \texttt{eliminate\_essential\_bc}(\mat{K},E,\mathtt{DIAG\_ZERO})
  \;=\; \mat{P}_F\,\mat{K}\,\mat{P}_F,
  \label{eq:pfproj}
\end{equation}
because left-multiplying by $\mat{P}_F$ zeroes the essential rows and
right-multiplying zeroes the essential columns. This makes the operation
\emph{linear} in $\mat{K}$, and linearity in turn gives a sharp statement of
separability: elimination distributes over the assembly fold's term-sum. If
$\mat{K}=\mat{K}_1+\mat{K}_2$ is the sum of two assembled term-contributions, then
under \texttt{DIAG\_ZERO},
\[
  \texttt{eliminate\_essential\_bc}(\mat{K}_1+\mat{K}_2,E)
  = \texttt{eliminate\_essential\_bc}(\mat{K}_1,E)
  + \texttt{eliminate\_essential\_bc}(\mat{K}_2,E).
\]
The \texttt{DIAG\_ONE} policy is the same projection plus the constant $\mat{I}_E$
on the essential block---an affine, not linear, map---so it distributes up to that
one added identity. Either way, elimination factors through the free-block
projection \emph{regardless of how the matrix was assembled}, which is the precise
content of ``separable post-composition.''

\subsection{Inhomogeneous Dirichlet data: lifting into the right-hand side}
\label{sec:lifting}

When the prescribed boundary values are zero---a grounded conductor, a PEC
wall---the forcing term $-\mat{K}_{FE}\vx_E$ in \cref{eq:reduced} vanishes,
because $\vx_E=\mathbf{0}$. The reduced system is just
$\mat{K}_{FF}\vx_F=\vb_F$ and there is nothing more to do.

But the interesting essential conditions are \emph{inhomogeneous}: a capacitor
plate held at $V\neq 0$, a port driven at a nonzero amplitude. Then $\vx_E\neq
\mathbf{0}$, and the term $-\mat{K}_{FE}\vx_E$ in \cref{eq:reduced} is a genuine,
nonzero load. This is the discrete face of the ``inhomogeneous essential''
condition we flagged back in the parallel-plate example of \cref{ch:weak}: a
problem with no volume source and no surface charge ($f=g=0$, an empty $\ell$)
still produces a nonzero right-hand side, driven entirely by the imposed
boundary values. The mechanism that produces it is called \emph{lifting}.

\begin{figure}[t]
\centering
\begin{tikzpicture}[scale=1.0, node distance=7mm]
  % a known boundary value
  \node[data, fill=simBgBlue, draw=simBlue, text width=20mm] (xe)
    {known boundary\\data $\vx_{\mathrm{ess}}$};
  % apply K
  \node[opbox, right=14mm of xe, text width=22mm] (apply)
    {apply operator\\$\mat{K}\cdot\vx_{\mathrm{ess}}$};
  % subtract from b
  \node[opbox, right=14mm of apply, text width=24mm] (sub)
    {subtract from RHS\\$\vb \leftarrow \vb - \mat{K}\vx_{\mathrm{ess}}$};
  \node[product, right=14mm of sub, text width=18mm, align=center] (out)
    {lifted RHS\\$\vb'$};
  \draw[flow] (xe) -- (apply);
  \draw[flow] (apply) -- (sub);
  \draw[flow] (sub) -- (out);
  % the b feeding in
  \node[data, above=8mm of sub, font=\scriptsize\itshape] (b) {original $\vb$};
  \draw[flow] (b) -- (sub);
  % annotation
  \node[font=\scriptsize\itshape, simRust, below=8mm of apply.south, text width=40mm, align=center]
    {the boundary value is ``lifted'' off the unknown and recast as a known forcing on the interior};
  \draw[refedge] (apply.south) -- ++(0,-3mm);
\end{tikzpicture}
\caption[Lifting an inhomogeneous Dirichlet condition into the RHS]{Lifting. A
known boundary value $\vx_{\mathrm{ess}}$ on the essential DoFs is not an unknown
to solve for---it is data. The operation \texttt{eliminate\_rhs} applies the
assembled operator to it, $\mat{K}\,\vx_{\mathrm{ess}}$, and \emph{subtracts the
result from the right-hand side}, $\vb\leftarrow\vb-\mat{K}\vx_{\mathrm{ess}}$.
The boundary value's influence on the interior equations is thereby moved from the
matrix (where it was a coupling) to the load vector (where it is a known source).
The essential rows of $\vb'$ are then pinned to the boundary data, matching the
identity diagonal that \texttt{eliminate\_essential\_bc} installed.}
\label{fig:lifting}
\end{figure}

Lifting is exactly the move from \cref{eq:reduced}, performed on the full index
space. Let $\vx_{\mathrm{ess}}$ be the true-DoF vector that holds the prescribed
boundary values on the essential DoFs (and zero elsewhere). The operation
\texttt{eliminate\_rhs} computes
\begin{equation}
  \vb' \;=\; \vb \;-\; \mat{K}\,\vx_{\mathrm{ess}},
  \label{eq:liftalg}
\end{equation}
then overwrites the essential rows of $\vb'$ with the boundary data itself
(matching the identity diagonal of the \texttt{DIAG\_ONE} policy). Read against
\cref{eq:reduced}: on a free row $i\in F$, the product $\mat{K}\vx_{\mathrm{ess}}$
contributes $\sum_{j\in E}\mat{K}_{ij}(\vx_{\mathrm{ess}})_j = (\mat{K}_{FE}\vx_E)_i$,
which is exactly the term being subtracted. On an essential row, the pin
overwrites whatever was there with the prescribed value, so the essential equation
reads $1\cdot x_i = (\vx_{\mathrm{ess}})_i$. The two operations together---
\texttt{eliminate\_essential\_bc} on the operator, \texttt{eliminate\_rhs} on the
right-hand side---realize the full Dirichlet treatment of \cref{eq:reduced}
without ever re-indexing the matrix.

Algebraically, \cref{eq:liftalg} is just one matrix--vector product
($\mat{K}\vx_{\mathrm{ess}}$) and one scaled vector addition ($\vb-(\cdot)$,
an \texttt{axpy} with coefficient $-1$). It carries no SPD or invertibility
requirement on $\mat{K}$; it simply moves a known quantity from one side of the
equation to the other. Notice, too, that lifting is a \emph{one-shot}
preparation, not a projector: apply it twice and you subtract
$\mat{K}\vx_{\mathrm{ess}}$ twice, which is wrong. It is performed exactly once,
on the way into the solve.

\begin{workedexample}{Lifting a nonzero Dirichlet value in 1-D}{lift1d}
We carry the whole essential-BC machinery through on a system small enough to
write out completely. Discretize the 1-D Poisson problem $-u''=0$ on $[0,1]$ with
the boundary conditions $u(0)=0$ and $u(1)=V$ (an inhomogeneous Dirichlet datum
at the right end), using a uniform grid with two interior nodes. The standard
finite-difference/finite-element stiffness on nodes $x_0,x_1,x_2,x_3$ (with
$x_0=0$, $x_3=1$, spacing $h=\tfrac13$) is, up to the factor $1/h$,
\[
  \mat{K} =
  \begin{bmatrix}
     1 & -1 &  0 &  0\\
    -1 &  2 & -1 &  0\\
     0 & -1 &  2 & -1\\
     0 &  0 & -1 &  1
  \end{bmatrix},
  \qquad \vb = \mathbf{0}.
\]
The essential set is $E=\{0,3\}$ (the two endpoints); the free set is $F=\{1,2\}$
(the two interior nodes). The prescribed boundary data is
$\vx_{\mathrm{ess}}=(0,\,?,\,?,\,V)^{\!\top}$, with the interior entries
irrelevant (they will be masked out).

\medskip\noindent\textbf{Reduced system, by hand.}
Writing $\vx=(x_0,x_1,x_2,x_3)$ with $x_0=0$, $x_3=V$ known, the two interior
equations (rows $1$ and $2$) are
\[
  -x_0 + 2x_1 - x_2 = 0, \qquad -x_1 + 2x_2 - x_3 = 0.
\]
Substituting $x_0=0$ and $x_3=V$ and moving the known terms to the right gives the
reduced free system \cref{eq:reduced}:
\[
  \begin{bmatrix} 2 & -1\\ -1 & 2\end{bmatrix}
  \begin{bmatrix} x_1\\ x_2\end{bmatrix}
  =
  \begin{bmatrix} 0\\ V\end{bmatrix}.
\]
The right-hand side entry $V$ in the second equation is precisely the lifted term
$-(\mat{K}_{FE}\vx_E)_2 = -(-1)\cdot V = V$. Solving:
$x_1 = V/3$, $x_2 = 2V/3$. Together with the endpoints this is the exact
linear potential $u(x)=Vx$ sampled at the four nodes
$(0,\,V/3,\,2V/3,\,V)$---as it must be, since the solution of $-u''=0$ with these
boundary values is the straight line from $0$ to $V$.

\medskip\noindent\textbf{The same thing via the two operations.}
Now reproduce the result through \texttt{eliminate\_rhs} and
\texttt{eliminate\_essential\_bc} on the full $4\times4$ system, without
re-indexing. First lift the RHS with $\vx_{\mathrm{ess}}=(0,0,0,V)^{\!\top}$:
\[
  \mat{K}\,\vx_{\mathrm{ess}} =
  \begin{bmatrix} 1&-1&0&0\\-1&2&-1&0\\0&-1&2&-1\\0&0&-1&1\end{bmatrix}
  \begin{bmatrix}0\\0\\0\\V\end{bmatrix}
  = \begin{bmatrix}0\\0\\-V\\V\end{bmatrix},
  \qquad
  \vb' = \vb - \mat{K}\vx_{\mathrm{ess}} =
  \begin{bmatrix}0\\0\\V\\-V\end{bmatrix}.
\]
On the free rows $\{1,2\}$ this gives $\vb'_1=0$, $\vb'_2=V$---exactly the
right-hand side of the reduced system above. (The essential rows $\{0,3\}$ are
then pinned to the boundary data $0$ and $V$.) Second, eliminate the operator with
\texttt{DIAG\_ONE}: zero rows/columns $0$ and $3$ and set their diagonal to $1$,
producing
\[
  \mat{K}_{\mathrm{BC}} =
  \begin{bmatrix} 1&0&0&0\\ 0&2&-1&0\\ 0&-1&2&0\\ 0&0&0&1\end{bmatrix}.
\]
Solving $\mat{K}_{\mathrm{BC}}\vx=\vb'$ with $\vb'=(0,0,V,-V)^{\!\top}$ after
pinning the essential rows to $(0,\dots,V)$ recovers the free block
$\bigl[\begin{smallmatrix}2&-1\\-1&2\end{smallmatrix}\bigr]
(x_1,x_2)^{\!\top}=(0,V)^{\!\top}$ and hence $x_1=V/3$, $x_2=2V/3$ again. The
matrix-level operations reproduce the by-hand reduction exactly, which is the
whole point: \texttt{eliminate\_essential\_bc} and \texttt{eliminate\_rhs} are a
re-indexing-free implementation of \cref{eq:reduced}.
\end{workedexample}

\subsection{Natural conditions, by contrast}
\label{sec:natural}

It is worth setting the essential machinery beside the natural conditions to see
how little the latter cost. Recall the boundary split of \cref{ch:weak}: Green's
identity produced a boundary integral $\oint_\bdry v\,(\varepsilon\Grad u)\cdot
\vn\,\mathrm{d}S$, and we disposed of it in two ways. On the Dirichlet part we
forced the test functions to vanish and the integral dropped---the essential
condition, which then had to be imposed on the discrete system by the elimination
we just built. On the Neumann part the integrand was a \emph{known} flux $g$, and
the surviving integral $\oint_{\Gamma_N} g\,v\,\mathrm{d}S$ moved to the
right-hand side as a source.

That natural-BC surface term is assembled \emph{during} assembly, not after. It is
just another weak-form term---a boundary integrator $(g,\text{value})$ on the
surface $\Gamma_N$, folded into $\vb$ exactly like a volume source---and the
assembly machinery of \cref{ch:assembly} handles it with no special case. A
\emph{homogeneous} Neumann condition ($g=0$) contributes nothing; it is enforced
by doing nothing at all. This is the asymmetry the pitfall in \cref{ch:weak}
warned about, now visible in implementation terms: the natural condition is a row
in the integrator list, while the essential condition is a post-assembly surgery
on rows and columns.

\begin{figure}[t]
\centering
\begin{tikzpicture}[scale=1.0]
  % essential side
  \begin{scope}[shift={(0,0)}]
    \node[font=\small\bfseries, simBlue] at (1.6,2.6) {essential (Dirichlet)};
    \node[stage, fill=simBgBlue, draw=simBlue, text width=30mm, align=center] (e1) at (1.6,1.7)
      {value $u=u_0$ prescribed};
    \node[opbox, text width=34mm, below=4mm of e1] (e2)
      {enforced on the \emph{discrete system}:\\zero rows/cols, lift RHS};
    \node[font=\scriptsize\itshape, simRust, below=3mm of e2] (e3)
      {\emph{after} assembly --- a post-composition};
    \draw[flow] (e1) -- (e2);
    \draw[flow] (e2) -- (e3);
  \end{scope}
  % divider
  \draw[simGray, densely dashed] (3.7,-0.7) -- (3.7,2.8);
  % natural side
  \begin{scope}[shift={(5.0,0)}]
    \node[font=\small\bfseries, simGreen] at (1.8,2.6) {natural (Neumann)};
    \node[stage, fill=simBgGreen, draw=simGreen, text width=30mm, align=center] (n1) at (1.8,1.7)
      {flux $(\varepsilon\Grad u)\cdot\vn = g$ prescribed};
    \node[opbox, text width=34mm, below=4mm of n1] (n2)
      {enters the \emph{weak form}:\\boundary integrator $\to \vb$};
    \node[font=\scriptsize\itshape, simRust, below=3mm of n2] (n3)
      {\emph{during} assembly --- $g=0$ costs nothing};
    \draw[flow] (n1) -- (n2);
    \draw[flow] (n2) -- (n3);
  \end{scope}
\end{tikzpicture}
\caption[Essential versus natural, in the discrete setting]{The two boundary-
condition types impose themselves at different stages. An \emph{essential}
condition (left) prescribes the field's \emph{value} and must be enforced on the
assembled system, after the fold, by zeroing rows/columns and lifting the data
into the right-hand side---the work of \cref{sec:essential}. A \emph{natural}
condition (right) prescribes a \emph{flux} and enters the weak form as a boundary
integrator folded into $\vb$ \emph{during} assembly; a homogeneous one ($g=0$)
costs nothing. ``Essential'' means ``you enforce it''; ``natural'' means ``the
weak form enforces it for you.''}
\label{fig:essvsnat}
\end{figure}

\section{The Helmholtz decomposition}
\label{sec:helmholtz}

We turn to the second finishing operation. To motivate it we need a classical
theorem about vector fields, and its discrete face.

\subsection{The continuous decomposition}

A foundational fact of vector calculus is that any reasonable vector field on a
domain can be split into two pieces of complementary character: a part with no
divergence and a part that is a pure gradient.

\begin{theorem}[Helmholtz decomposition]
\label{thm:helmholtz}
Let $\vF$ be a (sufficiently smooth) vector field on a bounded domain $\dom$ with
appropriate boundary conditions. Then $\vF$ decomposes uniquely as
\begin{equation}
  \vF \;=\; \vF_{\mathrm{df}} \;+\; \Grad\psi,
  \label{eq:helmholtz}
\end{equation}
where $\vF_{\mathrm{df}}$ is \emph{divergence-free} ($\Div\vF_{\mathrm{df}}=0$)
and $\Grad\psi$ is a \emph{gradient} (and hence curl-free, $\Curl\Grad\psi=0$).
The two parts are orthogonal in the $\Ltwo$ inner product.
\end{theorem}

The two pieces are sometimes called the \emph{solenoidal} (divergence-free) and
\emph{irrotational} (gradient) components. Physically they are the field's
``circulating'' part and its ``spreading'' part: a magnetic field is purely
solenoidal ($\Div\vB=0$ always), while an electrostatic field is purely
irrotational ($\vE=-\Grad V$). A general field mixes the two, and
\cref{thm:helmholtz} says the mixture separates cleanly and uniquely.
\Cref{fig:helmholtz} shows the decomposition on a small field.

The orthogonality claim is worth a moment, because it is the property the discrete
projector inherits and the reason the decomposition is \emph{unique}. The two
parts are $\Ltwo$-orthogonal:
\begin{equation}
  \inner{\vF_{\mathrm{df}}}{\Grad\psi}_{\Ltwo(\dom)}
  = \int_\dom \vF_{\mathrm{df}}\cdot\Grad\psi \dV = 0.
  \label{eq:helmorth}
\end{equation}
The reason is Green's first identity (\cref{ch:weak}) run backwards. Integrate by
parts: $\int_\dom \vF_{\mathrm{df}}\cdot\Grad\psi
= -\int_\dom (\Div\vF_{\mathrm{df}})\,\psi
+ \oint_\bdry \psi\,(\vF_{\mathrm{df}}\cdot\vn)\,\mathrm{d}S$. The volume term
vanishes because $\vF_{\mathrm{df}}$ is divergence-free ($\Div\vF_{\mathrm{df}}=0$),
and the boundary term vanishes under the boundary conditions that accompany the
decomposition (either $\psi=0$ or $\vF_{\mathrm{df}}\cdot\vn=0$ on $\bdry$). So a
divergence-free field is $\Ltwo$-orthogonal to every gradient---which is precisely
what makes ``the gradient part'' well-defined: it is the orthogonal projection of
$\vF$ onto the gradient subspace, and $\vF_{\mathrm{df}}$ is whatever remains.
Uniqueness follows immediately: if $\vF=\vF_{\mathrm{df}}+\Grad\psi
=\vF'_{\mathrm{df}}+\Grad\psi'$ were two decompositions, their difference
$\vF_{\mathrm{df}}-\vF'_{\mathrm{df}}=\Grad(\psi'-\psi)$ would be both
divergence-free and a gradient, hence orthogonal to itself, hence zero. The
discrete projector of \cref{sec:divfree} reproduces this orthogonality---but in a
\emph{weighted} inner product, the $\mat{M}$-inner-product, because the discrete
mass operator $\mat{M}$ encodes the material coefficient $\varepsilon$.

\begin{figure}[t]
\centering
\begin{tikzpicture}[scale=1.0,
    vec/.style={-{Stealth[length=2mm]}, thick}]
  % --- the full field F (left) ---
  \begin{scope}[shift={(0,0)}]
    \node[font=\small\bfseries] at (1.5,3.5) {$\vF$ (mixed)};
    \draw[simGray!50] (0,0) rectangle (3,3);
    \foreach \x in {0.5,1.5,2.5}{
      \foreach \y in {0.5,1.5,2.5}{
        % a swirling-plus-spreading field
        \pgfmathsetmacro{\dx}{0.32*cos(120*\x+60*\y) + 0.20*(\x-1.5)}
        \pgfmathsetmacro{\dy}{0.32*sin(120*\x+60*\y) + 0.20*(\y-1.5)}
        \draw[vec,simInk] (\x,\y) -- ++(\dx,\dy);}}
  \end{scope}
  % equals
  \node[font=\Large] at (3.7,1.5) {$=$};
  % --- divergence-free part (middle) ---
  \begin{scope}[shift={(4.4,0)}]
    \node[font=\small\bfseries, simBlue] at (1.5,3.5) {$\vF_{\mathrm{df}}$ (curl)};
    \draw[simGray!50] (0,0) rectangle (3,3);
    \foreach \x in {0.5,1.5,2.5}{
      \foreach \y in {0.5,1.5,2.5}{
        \pgfmathsetmacro{\dx}{0.32*cos(120*\x+60*\y)}
        \pgfmathsetmacro{\dy}{0.32*sin(120*\x+60*\y)}
        \draw[vec,simBlue] (\x,\y) -- ++(\dx,\dy);}}
  \end{scope}
  % plus
  \node[font=\Large] at (8.1,1.5) {$+$};
  % --- gradient part (right) ---
  \begin{scope}[shift={(8.8,0)}]
    \node[font=\small\bfseries, simRust] at (1.5,3.5) {$\Grad\psi$ (spread)};
    \draw[simGray!50] (0,0) rectangle (3,3);
    \foreach \x in {0.5,1.5,2.5}{
      \foreach \y in {0.5,1.5,2.5}{
        \pgfmathsetmacro{\dx}{0.20*(\x-1.5)}
        \pgfmathsetmacro{\dy}{0.20*(\y-1.5)}
        \draw[vec,simRust] (\x,\y) -- ++(\dx,\dy);}}
    % a contour of psi
    \draw[simRust!50,densely dotted] (1.5,1.5) circle (0.5);
    \draw[simRust!50,densely dotted] (1.5,1.5) circle (1.0);
  \end{scope}
\end{tikzpicture}
\caption[The Helmholtz decomposition of a vector field]{The Helmholtz
decomposition \cref{eq:helmholtz} splits a vector field $\vF$ (left, mixed) into a
\emph{divergence-free} part $\vF_{\mathrm{df}}$ (middle, purely circulating---it
swirls but neither spreads from nor converges to any point) and a \emph{gradient}
part $\Grad\psi$ (right, purely spreading---it points along the steepest ascent of
a scalar potential $\psi$, whose level sets are the dotted contours). The two
parts are $\Ltwo$-orthogonal. The divergence-free part is what the physics of a
magnetic or resonant field demands; the gradient part is the unphysical
contamination a singular operator can introduce.}
\label{fig:helmholtz}
\end{figure}

\subsection{Its de Rham meaning}

The decomposition is not an isolated trick; it is a face of the de Rham complex of
\cref{ch:derham}. Recall the ladder of spaces and operators
\[
  \Hone \xrightarrow{\ \Grad\ } \Hcurl \xrightarrow{\ \Curl\ }
  \Hdiv \xrightarrow{\ \Div\ } \Ltwo,
\]
with the defining \emph{exactness} property that the range of each operator is the
kernel of the next: $\Curl\Grad = 0$ and $\Div\Curl = 0$, and (on a
topologically simple domain) the inclusions are equalities. Look at the middle
slot, $\Hcurl$. The gradient part of \cref{eq:helmholtz} lives in
$\mathrm{Range}(\Grad)$, the image of the gradient operator coming \emph{into}
$\Hcurl$ from $\Hone$. Exactness says this image is exactly $\mathrm{Ker}(\Curl)$,
the curl-free fields. So the Helmholtz split of an $\Hcurl$ field is precisely its
decomposition into
\begin{equation}
  \Hcurl \;=\; \underbrace{\mathrm{Range}(\Grad)}_{\text{gradients}\,=\,\mathrm{Ker}(\Curl)}
  \;\oplus\;
  \underbrace{(\text{divergence-free complement})}_{\vF_{\mathrm{df}}},
  \label{eq:hcurlsplit}
\end{equation}
the gradient part being the kernel of the curl. This is the structural reason the
decomposition matters for us: the operators we solve live on $\Hcurl$, and their
null space \emph{is} the gradient subspace of \cref{eq:hcurlsplit}.

\begin{keyidea}
The Helmholtz decomposition \cref{eq:helmholtz} is the de Rham complex read at the
$\Hcurl$ slot. The gradient part of an $\Hcurl$ field is its component in
$\mathrm{Range}(\Grad)$, which by exactness of the complex
(\cref{ch:derham}) is exactly $\mathrm{Ker}(\Curl)$---the fields the curl
annihilates. Removing the gradient part is therefore the same as removing the
curl operator's null space. The discrete projector of the next section is a
matrix realization of this removal.
\end{keyidea}

\section{The discrete divergence-free projector}
\label{sec:divfree}

We now build the operator that removes the gradient part from a discrete field.
It is the workhorse $\mathtt{divfree\_project}$, and it is, like essential
elimination, a separable post-composition---but a richer one, because it carries
an inner linear solve.

\subsection{The projector and its apply}

Let $\mat{G}$ be the discrete gradient operator (the matrix that takes the
$\Hone$ DoFs of a scalar potential to the $\Hcurl$ DoFs of its gradient---the de
Rham interpolator of \cref{ch:derham}), and let $\mat{M}$ be the
$\varepsilon$-weighted $\Hone$ mass operator (symmetric positive definite). The
discrete divergence-free projector is
\begin{equation}
  \boxed{\;
  \mat{P} \;=\; \mat{I} \;-\; \mat{G}\,\bigl(\mat{G}^{\!\top}\mat{M}\,\mat{G}\bigr)^{-1}\,\mat{G}^{\!\top}\mat{M}.
  \;}
  \label{eq:projector}
\end{equation}
This is an $\mat{M}$-orthogonal projection onto the divergence-free subspace. Read
it from right to left as a recipe: $\mat{G}^{\!\top}\mat{M}$ measures the (weak)
divergence of the input field; $(\mat{G}^{\!\top}\mat{M}\mat{G})^{-1}$ solves a
Poisson-type system for a scalar potential $\psi$; $\mat{G}$ takes the gradient of
that potential, materializing the gradient component; and $\mat{I}-(\cdots)$
subtracts it off, leaving the divergence-free remainder. The middle factor
$(\mat{G}^{\!\top}\mat{M}\mat{G})^{-1}$ is the $\mat{M}$-weighted normal-equations
inverse: $\psi$ is chosen so that $\mat{G}\psi$ is the $\mat{M}$-orthogonal
projection of the input onto the gradient subspace.

\begin{notationbox}
The triple product $\mat{G}^{\!\top}\mat{M}\mat{G}$ is the Poisson-type stiffness
on the scalar ($\Hone$) space---it is what you get by sandwiching the mass
operator between the gradient and its transpose. In Palace it is never formed as a
matrix: the inner solve is posed against $\mat{M}$ directly, with $\mat{G}^{\!\top}$
applied on the right-hand side (as a weak-divergence operator) and $\mat{G}$
applied to the solution. The boxed form \cref{eq:projector} is the algebra; the
implementation realizes it factor by factor, which is why the four-step apply
below mentions $\mat{M}$, $\mat{G}$, and a weak-divergence operator separately
rather than the triple product.
\end{notationbox}

In practice we never form $\mat{P}$ as a matrix and we rarely apply the full
$\mat{P}$; we apply it to a field $\vy$ to extract the divergence-free part
$\vy'=\mat{P}\vy$. The apply runs in four steps, shown as a pipeline in
\cref{fig:projapply}:

\begin{enumerate}\setlength\itemsep{2pt}
  \item \textbf{Weak divergence.} Compute $\vr = \mat{G}^{\!\top}\mat{M}\,\vy$, an
  $\Hone$-side residual measuring the divergence of $\vy$. (In Palace this is one
  application of a weak-divergence operator, the $\varepsilon$-weighted
  Nédélec-to-$\Hone$ map; its sign convention is load-bearing and discussed below.)
  \item \textbf{Essential-BC zeroing.} Zero $\vr$ on the essential boundary DoFs of
  the scalar space, $\vr\leftarrow \mat{Z}_{\mathrm{bdr}}\vr$. This pins the inner
  Poisson problem and removes the pure-Neumann null space (a constant shift in
  $\psi$ that would otherwise leave the solve singular).
  \item \textbf{Inner Poisson solve.} Solve $\mat{M}\psi = \vr$ for the scalar
  potential $\psi$. This is itself an iterative \texttt{ksp\_solve}
  (\cref{ch:krylov})---a conjugate-gradient solve against the SPD operator
  $\mat{M}$, preconditioned. The projector therefore \emph{contains a solver}: it
  is a constructed operator carrying an inner solver as a sub-component
  (\cref{ch:operators}).
  \item \textbf{Gradient correction.} Form the gradient $\mat{G}\psi$ and subtract
  it: $\vy' = \vy - \mat{G}\psi$ (an \texttt{apply\_linop} followed by an
  \texttt{axpy}). The result $\vy'$ is the divergence-free component.
\end{enumerate}

\begin{figure}[t]
\centering
\begin{tikzpicture}[node distance=8mm]
  \node[data, text width=12mm, align=center] (y) {field $\vy$};
  \node[opbox, right=9mm of y, text width=20mm] (wd)
    {\textbf{1.}\,weak div\\$\mat{G}^{\!\top}\mat{M}\vy$};
  \node[opbox, right=8mm of wd, text width=18mm] (z)
    {\textbf{2.}\,zero BC\\$\mat{Z}_{\mathrm{bdr}}$};
  \node[extern, right=8mm of z, text width=22mm] (solve)
    {\textbf{3.}\,inner solve\\$\mat{M}\psi=\vr$\\[1pt]\scriptsize \texttt{ksp\_solve}};
  \node[opbox, below=10mm of solve, text width=20mm] (grad)
    {\textbf{4.}\,gradient\\$\mat{G}\psi$};
  \node[product, left=8mm of grad, text width=22mm, align=center] (sub)
    {subtract\\$\vy'=\vy-\mat{G}\psi$};
  \node[data, left=8mm of sub, text width=18mm, align=center] (out)
    {divergence-free $\vy'$};
  \draw[flow] (y) -- (wd);
  \draw[flow] (wd) -- (z);
  \draw[flow] (z) -- (solve);
  \draw[flow] (solve) -- (grad);
  \draw[flow] (grad) -- (sub);
  \draw[flow] (sub) -- (out);
  % the original y feeding the subtract
  \draw[flow, simGray, densely dashed] (y.south) |- (sub.north);
  \node[font=\scriptsize\itshape, simViolet, above=1mm of solve.north]
    {the inner Krylov solve};
\end{tikzpicture}
\caption[The four-step divergence-free projector apply]{The
$\mathtt{divfree\_project}$ apply as a pipeline. Given a field $\vy$, step~1
computes its weak divergence $\mat{G}^{\!\top}\mat{M}\vy$; step~2 zeroes that
residual on the essential boundary DoFs (pinning the inner problem); step~3
\emph{solves} a Poisson-type system $\mat{M}\psi=\vr$ for the scalar potential
$\psi$---an inner \texttt{ksp\_solve} (\cref{ch:krylov}), the dashed/violet box,
which makes the projector a constructed operator carrying its own solver; step~4
forms the gradient $\mat{G}\psi$ and subtracts it from the original $\vy$,
leaving the divergence-free component $\vy'$. Steps 1, 2, 4 are cheap whole-field
operations; the cost lives in the inner solve of step~3.}
\label{fig:projapply}
\end{figure}

The defining property is that the output is annihilated by the weak divergence:
$\mat{G}^{\!\top}\mat{M}\,\vy' = 0$. That is the discrete statement of
``divergence-free.'' And the part removed, $\mat{G}\psi$, is a discrete
gradient---the kernel of the discrete curl, by the de Rham exactness of
\cref{ch:derham}. So the projector strips exactly the curl operator's null space,
which is what \cref{eq:hcurlsplit} promised.

\begin{workedexample}{Walking the four-step apply on numbers}{fourstep}
To make the pipeline of \cref{fig:projapply} concrete, run all four steps on a
minimal instance with no essential boundary (so step~2 is the identity and we can
watch steps 1, 3, 4 do the work). Take a field space of dimension $2$ and a scalar
space of dimension $1$, with
\[
  \mat{M} = \mat{I}_2,
  \qquad
  \mat{G} = \begin{bmatrix}1\\1\end{bmatrix},
  \qquad
  \mat{G}^{\!\top}\mat{M} = \begin{bmatrix}1 & 1\end{bmatrix},
\]
the same data as \cref{wex:projcheck}. Project the field
$\vy=(3,1)^{\!\top}$---a mix of a circulating part along $(1,-1)$ and a gradient
part along $(1,1)$.

\medskip\noindent\textbf{Step 1 --- weak divergence.}
$\vr = \mat{G}^{\!\top}\mat{M}\,\vy = \begin{bmatrix}1&1\end{bmatrix}
\begin{bmatrix}3\\1\end{bmatrix} = 4$. The field has nonzero weak divergence, so it
is not already divergence-free---there is a gradient part to remove.

\medskip\noindent\textbf{Step 2 --- essential-BC zeroing.}
With no essential boundary, $\mat{Z}_{\mathrm{bdr}}=\mat{I}$ and $\vr$ passes
through unchanged: $\vr = 4$.

\medskip\noindent\textbf{Step 3 --- inner Poisson solve.}
Solve $(\mat{G}^{\!\top}\mat{M}\mat{G})\,\psi = \vr$, i.e.\ $2\psi = 4$, giving
$\psi = 2$. (In the implementation this is the inner \texttt{ksp\_solve} against
$\mat{M}$; here the scalar system is $1\times1$ and solves in closed form.) The
scalar potential whose gradient is the irrotational part of $\vy$ is $\psi=2$.

\medskip\noindent\textbf{Step 4 --- gradient correction.}
Form $\mat{G}\psi = \begin{bmatrix}1\\1\end{bmatrix}\cdot 2 = (2,2)^{\!\top}$ and
subtract: $\vy' = \vy - \mat{G}\psi = (3,1)^{\!\top} - (2,2)^{\!\top} =
(1,-1)^{\!\top}$.

\medskip\noindent\textbf{Check.}
The output $\vy'=(1,-1)^{\!\top}$ is exactly the antidiagonal (circulating)
component of $\vy$. Verify it is divergence-free:
$\mat{G}^{\!\top}\mat{M}\,\vy' = \begin{bmatrix}1&1\end{bmatrix}
\begin{bmatrix}1\\-1\end{bmatrix} = 0$, as the defining property demands. And the
removed part $\mat{G}\psi=(2,2)^{\!\top}$ is a pure gradient, along $(1,1)$. This
matches the closed-form projector of \cref{wex:projcheck}:
$\mat{P}\,\vy=\tfrac12\bigl[\begin{smallmatrix}1&-1\\-1&1\end{smallmatrix}\bigr]
\bigl[\begin{smallmatrix}3\\1\end{smallmatrix}\bigr]=\tfrac12(2,-2)^{\!\top}
=(1,-1)^{\!\top}$---the four-step apply and the boxed formula agree, as they must.
The point of running the steps separately is that the implementation never forms
$\mat{P}$: it walks exactly steps 1--4, and the expensive one is step~3, the inner
solve that here was a trivial $2\psi=4$ but in production is a full
preconditioned CG solve on the scalar space.
\end{workedexample}

\begin{notationbox}
\emph{A sign subtlety, recorded.} In Palace the weak-divergence operator carries a
built-in $-1$ (its bilinear form is $a(\vu,v)=-(\varepsilon\vu,\Grad v)$, so it is
the negated gradient transpose, $\mat{G}^{\!\top}$ up to that sign). The
implementation then \emph{adds} the correction, $\vy'\leftarrow\vy+\mat{G}\psi$,
and the two minus signs combine so that the net effect still removes the gradient
part. The algebra of \cref{eq:projector} is written with the conventional sign
($\vy-\mat{G}\psi$); the source's two-negatives bookkeeping is equivalent. We note
it because a single flipped sign here would invert the correction and the
projector would \emph{amplify} the gradient part instead of removing it.
\end{notationbox}

\subsection{It is a constructed operator with an inner solve}

Step 3 deserves emphasis. The projector is not a fixed sparse matrix; it is a
\emph{constructed operator}---a structured value assembled once at solver setup,
carrying the operators $\mat{M}$, $\mat{G}$, the weak-divergence operator, the
essential-boundary DoF set, and, crucially, an \emph{inner linear solver} bound to
$\mat{M}$. Applying the projector invokes that inner solver. This makes
$\mathtt{divfree\_project}$ the first operator we have met whose own definition
contains a solve-to-tolerance, and it foreshadows the constructed-operator
machinery of \cref{ch:operators}: operators that carry other operators (and other
solvers) as sub-components, applied opaquely. The inner conjugate-gradient
iteration is sequential---a recurrence that does not parallelize across its
steps---but that sequentiality is \emph{interior} to the inner solve; the
projector itself is a fixed four-step straight-line composition around it.

The construction-versus-apply split is the same one we saw for assembly, lifted one
level. There is a \emph{setup} phase, run once: build $\mat{M}$ and the
weak-divergence operator by assembly (themselves folds over weak-form terms on the
scalar space), interpolate the discrete gradient $\mat{G}$ from the de Rham
structure, compute the essential-boundary DoF set, and configure the inner solver
with its preconditioner, tolerance, and iteration cap. None of this depends on the
field being projected; it is bound into the constructed value $\mat{P}$ at solver
setup. Then there is the \emph{apply} phase, run once per field (and, in an
eigensolve, once per iterate): the four cheap-plus-one-expensive steps of
\cref{fig:projapply}. Keeping the two phases separate is what makes the per-iterate
projection affordable---the operators are assembled and the solver is set up exactly
once, and each projection reuses them.

The inner solve is also where the projector's cost concentrates, and where its one
genuine subtlety lives. The idempotence law $\mat{P}^2=\mat{P}$ of
\cref{wex:projcheck} rested on the middle factor being the \emph{exact} inverse
$(\mat{G}^{\!\top}\mat{M}\mat{G})^{-1}$. In practice the inner CG solve is only
approximate, converged to a tolerance, so $\mat{P}$ is an \emph{approximate}
projector: $\mat{P}^2\approx\mat{P}$, and $\mat{G}^{\!\top}\mat{M}\,\vy'\approx 0$,
both up to that tolerance. This is why the inner solve must be converged
tightly---loose tolerances leave a residual gradient component that the eigensolver
will happily amplify back into a spurious mode. The mass operator $\mat{M}$ is SPD,
so CG is the natural choice and a good preconditioner (an algebraic multigrid
$\mat{V}$-cycle on $\mat{M}$) makes the solve cheap; but cheap and \emph{loose} are
different knobs, and only the first may be turned.

\begin{remark}[The projector is a separable post-composition too]
Like essential elimination, the divergence-free projection consumes its inputs
opaquely. The four-step apply treats $\mat{M}$, $\mat{G}$, and the weak-divergence
operator as given operators and the field $\vy$ as a given vector; it never inspects
how any of them were assembled, nor which weak-form terms produced the curl--curl
operator it is guarding. So although the projector is structurally richer than
\texttt{eliminate\_essential\_bc}---it carries an inner solver where elimination
carries only an index set---it occupies the same architectural slot: a finishing
post-composition that acts on assembled objects and composes cleanly behind the
assembly fold.
\end{remark}

\subsection{Verifying the projector properties}

A projector must satisfy $\mat{P}^2=\mat{P}$ (applying it twice is the same as
once) and must annihilate the subspace it projects out (here, the gradients:
$\mat{P}\,\mat{G}\psi = 0$ for any $\psi$). Both follow from the algebraic form
\cref{eq:projector}, and both are worth seeing concretely.

\begin{workedexample}{$\mat{P}$ is a projector and kills gradients}{projcheck}
We verify the two projector laws, first symbolically from \cref{eq:projector},
then on a tiny numerical instance.

\medskip\noindent\textbf{Symbolic: idempotence.}
Abbreviate the middle inverse as $\mat{S}=(\mat{G}^{\!\top}\mat{M}\mat{G})^{-1}$,
so that $\mat{P}=\mat{I}-\mat{G}\,\mat{S}\,\mat{G}^{\!\top}\mat{M}$. Compute
$\mat{P}^2$:
\[
  \mat{P}^2 = \bigl(\mat{I}-\mat{G}\mat{S}\mat{G}^{\!\top}\mat{M}\bigr)
              \bigl(\mat{I}-\mat{G}\mat{S}\mat{G}^{\!\top}\mat{M}\bigr)
            = \mat{I} - 2\,\mat{G}\mat{S}\mat{G}^{\!\top}\mat{M}
              + \mat{G}\mat{S}\underbrace{\mat{G}^{\!\top}\mat{M}\mat{G}}_{=\,\mat{S}^{-1}}\mat{S}\,\mat{G}^{\!\top}\mat{M}.
\]
In the last term the inner $\mat{G}^{\!\top}\mat{M}\mat{G}$ is exactly
$\mat{S}^{-1}$, so $\mat{S}\,\mat{S}^{-1}\,\mat{S}=\mat{S}$ and the term collapses
to $\mat{G}\mat{S}\mat{G}^{\!\top}\mat{M}$. Hence
\[
  \mat{P}^2 = \mat{I} - 2\,\mat{G}\mat{S}\mat{G}^{\!\top}\mat{M}
              + \mat{G}\mat{S}\mat{G}^{\!\top}\mat{M}
            = \mat{I} - \mat{G}\mat{S}\mat{G}^{\!\top}\mat{M} = \mat{P}.
\]
The cancellation of the middle inverse against the triple product is the entire
content of the projector law; it is why the recipe must use the \emph{exact}
inverse $\mat{S}$ (and why the inner solve must be converged tightly).

\medskip\noindent\textbf{Symbolic: gradients are killed.}
Apply $\mat{P}$ to a gradient $\mat{G}\psi$:
\[
  \mat{P}\,\mat{G}\psi
  = \mat{G}\psi - \mat{G}\mat{S}\mat{G}^{\!\top}\mat{M}\mat{G}\psi
  = \mat{G}\psi - \mat{G}\mat{S}\,\mat{S}^{-1}\psi
  = \mat{G}\psi - \mat{G}\psi = 0,
\]
again using $\mat{G}^{\!\top}\mat{M}\mat{G}=\mat{S}^{-1}$. So every gradient is in
the kernel of $\mat{P}$: $\mathrm{Ker}(\mat{P})=\mathrm{Range}(\mat{G})$, exactly
the curl null space of \cref{eq:hcurlsplit}.

\medskip\noindent\textbf{Numerical instance.}
Take the smallest nontrivial case: $\mat{M}=\mat{I}_2$ (the $2\times2$ identity,
so the $\mat{M}$-inner-product is ordinary $\Ltwo$), and a discrete gradient with
a single scalar DoF, $\mat{G}=\bigl[\begin{smallmatrix}1\\1\end{smallmatrix}\bigr]$
(the gradient subspace is the diagonal line spanned by $(1,1)$). Then
\[
  \mat{G}^{\!\top}\mat{M}\mat{G} = \begin{bmatrix}1&1\end{bmatrix}\begin{bmatrix}1\\1\end{bmatrix} = 2,
  \qquad \mat{S} = \tfrac12,
  \qquad
  \mat{G}\mat{S}\mat{G}^{\!\top}\mat{M}
  = \tfrac12\begin{bmatrix}1\\1\end{bmatrix}\begin{bmatrix}1&1\end{bmatrix}
  = \tfrac12\begin{bmatrix}1&1\\1&1\end{bmatrix}.
\]
Therefore
\[
  \mat{P} = \mat{I} - \tfrac12\begin{bmatrix}1&1\\1&1\end{bmatrix}
          = \tfrac12\begin{bmatrix}1&-1\\-1&1\end{bmatrix}.
\]
Check idempotence directly:
$\mat{P}^2=\tfrac14\bigl[\begin{smallmatrix}2&-2\\-2&2\end{smallmatrix}\bigr]
=\tfrac12\bigl[\begin{smallmatrix}1&-1\\-1&1\end{smallmatrix}\bigr]=\mat{P}$. Check
that gradients die: for any scalar $\psi$,
$\mat{P}\,\mat{G}\psi = \tfrac12\bigl[\begin{smallmatrix}1&-1\\-1&1\end{smallmatrix}\bigr]
\bigl[\begin{smallmatrix}\psi\\\psi\end{smallmatrix}\bigr]
= \tfrac12\bigl[\begin{smallmatrix}\psi-\psi\\-\psi+\psi\end{smallmatrix}\bigr]=\mathbf{0}$.
And the projection of a general field is its divergence-free (here, orthogonal-to-
diagonal) part: $\mat{P}\,(a,b)^{\!\top}=\tfrac12(a-b,\,b-a)^{\!\top}$, which is
the component of $(a,b)$ along the antidiagonal $(1,-1)$---perpendicular, in the
$\mat{M}=\mat{I}$ inner product, to the gradient direction $(1,1)$. The projector
keeps the circulating part and discards the spreading part, exactly as
\cref{fig:helmholtz} drew it.
\end{workedexample}

\section{Why the projection is needed: the curl--curl null space}
\label{sec:nullspace}

We can now say precisely why a simulator bothers with all of this. The curl--curl
operator of the magnetostatic and eigenmode problems (\cref{ch:reductions-pde})
is, by construction, singular: it has an enormous null space.

Recall the curl--curl stiffness $\mat{K}$ assembled from the bilinear form
$\int_\dom \nu\,(\Curl\vu)\cdot(\Curl\vv)$ of \cref{ch:weak}. Any gradient field
$\vu=\Grad\phi$ has zero curl ($\Curl\Grad\phi=0$, de Rham exactness,
\cref{ch:derham}), so it contributes nothing to the form: $\mat{K}\,\Grad\phi=0$.
The entire gradient subspace---which, on a fine mesh, is a substantial fraction of
all the DoFs---lies in the null space of $\mat{K}$. The operator is hugely
rank-deficient (\cref{fig:nullspace}).

\begin{figure}[t]
\centering
\begin{tikzpicture}[scale=1.0]
  % the whole Hcurl space as a big ellipse
  \draw[simInk,thick] (0,0) ellipse (3.6 and 2.2);
  \node[font=\small,simInk] at (0,2.55) {$\Hcurl$ (all field DoFs)};
  % the gradient null-space as a big shaded sub-blob
  \begin{scope}
    \clip (0,0) ellipse (3.6 and 2.2);
    \fill[simBgRust] (-3.6,-2.2) rectangle (0.4,2.2);
  \end{scope}
  \draw[simRust,thick] (0.4,-2.16) -- (0.4,2.16);
  \node[font=\small,simRust,align=center] at (-1.7,0)
    {$\mathrm{Ker}(\mat{K})$\\$=\mathrm{Range}(\Grad)$\\[1pt]\scriptsize gradient null space\\\scriptsize(spurious zero modes)};
  % the physical divergence-free subspace
  \node[font=\small,simBlue,align=center] at (1.9,0)
    {divergence-\\free part\\[1pt]\scriptsize the physics};
  % an eigensolver "drifting" into the null space
  \draw[backflow] (2.6,1.3) .. controls (1.0,1.6) and (-0.8,1.4) .. (-2.2,0.9);
  \node[font=\scriptsize\itshape,simRust] at (0.2,1.75) {solver drifts in};
  % the projector pushing back out
  \draw[depedge] (-2.0,-0.9) .. controls (-0.6,-1.5) and (1.0,-1.4) .. (2.4,-0.9);
  \node[font=\scriptsize\itshape,simBlue] at (0.2,-1.75) {$\mat{P}$ projects out};
\end{tikzpicture}
\caption[The curl--curl gradient null space]{Why the projection is needed. The
curl--curl operator $\mat{K}$ annihilates every gradient field
($\mat{K}\Grad\phi=0$, de Rham exactness), so its null space is the entire
gradient subspace $\mathrm{Range}(\Grad)$ (rust)---a large fraction of all the
field DoFs. These are \emph{spurious zero modes}: eigensolutions with eigenvalue
zero that carry no physics. An eigensolver or iterative solver, left unguarded,
drifts into this subspace (dashed) and returns garbage or stalls. The
divergence-free projector $\mat{P}$ pushes every iterate back onto the physical
divergence-free subspace (blue), so the solve sees only the well-posed part of the
operator.}
\label{fig:nullspace}
\end{figure}

For a solver this is poison in two ways. An \emph{eigensolver}
(\cref{ch:eigenmodes}) looking for the resonant modes
$\mat{K}\vx=\lambda\mat{M}\vx$ will find, in addition to the physical modes, a
flood of spurious modes with eigenvalue $\lambda\approx 0$---the gradient null
space, masquerading as zero-frequency ``resonances'' that are nothing of the kind.
They swamp the spectrum near zero and bury the physical low-frequency modes the
analysis actually wants. An \emph{iterative linear solver} fares no better: the
null space makes the operator singular (or, after a tiny regularization,
catastrophically ill-conditioned), and a Krylov method either stalls or wanders
into the null space and returns a contaminated answer.

\begin{pitfall}
The curl--curl operator's gradient null space is not a numerical artifact to be
regularized away with a small diagonal shift---it is a structural feature of the
physics and the de Rham complex. A shift would perturb the physical spectrum and
still leave near-null modes to confuse the solver. The correct cure is to
\emph{project the null space out} with $\mat{P}$, restricting the solve to the
divergence-free subspace where the operator is genuinely well-posed. Eigensolvers
in particular must project \emph{every} iterate (the initial vector and each
subsequent candidate), or the spurious modes creep back in.
\end{pitfall}

The remedy is the projector $\mat{P}$ of \cref{sec:divfree}. By construction its
range is the divergence-free subspace and its kernel is exactly the gradient null
space $\mathrm{Range}(\mat{G})=\mathrm{Ker}(\mat{K})$. Applying $\mat{P}$ to a
field removes precisely the troublesome component and leaves the physical part. In
the eigenmode analysis (\cref{ch:eigenmodes}) the projector is threaded into the
eigenvalue iteration: the initial vector is projected once, and every candidate
eigenvector is projected at each step, so the iteration is confined to the
divergence-free subspace throughout and never sees the spurious modes
(\cref{fig:eigproj}). The projector is the guardrail that keeps the eigensolver on
the physical manifold.

\begin{figure}[t]
\centering
\begin{tikzpicture}[node distance=8mm]
  \node[data, text width=16mm, align=center] (v0) {start vector $\vv_0$};
  \node[opbox, right=8mm of v0, text width=14mm, fill=simBgGreen, draw=simGreen] (p0)
    {$\mat{P}$ project};
  \node[stage, right=8mm of p0, text width=24mm, align=center] (step)
    {eigen step\\[1pt]\footnotesize (\cref{ch:eigenmodes})};
  \node[opbox, right=8mm of step, text width=14mm, fill=simBgGreen, draw=simGreen] (pk)
    {$\mat{P}$ project};
  \node[product, right=8mm of pk, text width=16mm, align=center] (conv)
    {converged?};
  \draw[flow] (v0) -- (p0);
  \draw[flow] (p0) -- (step);
  \draw[flow] (step) -- (pk);
  \draw[flow] (pk) -- (conv);
  % loop back
  \draw[backflow] (conv.south) -- ++(0,-7mm) -| node[below, pos=0.25, font=\scriptsize]{no: next candidate} (step.south);
  % output
  \node[data, right=8mm of conv, text width=16mm, align=center, fill=simBgBlue, draw=simBlue] (out) {physical mode};
  \draw[flow] (conv) -- node[above, font=\scriptsize]{yes} (out);
  \node[font=\scriptsize\itshape, simGreen, above=1mm of pk.north] {every iterate stays divergence-free};
\end{tikzpicture}
\caption[The projector inside the eigenvalue loop]{How the divergence-free
projector guards an eigensolve. The initial vector $\vv_0$ is projected once, and
\emph{every} candidate produced by the eigen step (\cref{ch:eigenmodes}) is
projected again with $\mat{P}$ before the convergence test and before it feeds the
next iteration (green boxes). Because $\mathrm{Ker}(\mat{P})$ is exactly the
gradient null space, the iteration is confined to the divergence-free subspace
throughout: the spurious zero modes of \cref{fig:nullspace} can never re-enter, and
the solver converges to the physical resonances alone. Projecting only the initial
vector is not enough---rounding and the iteration itself would let the null space
leak back in, which is why the projection sits \emph{inside} the loop.}
\label{fig:eigproj}
\end{figure}

\section{Closing Part II}

With boundary conditions imposed and fields projected onto their physical
subspace, the discretized problem is complete. Let us take stock of what Part~II
has built. We began (\cref{ch:weak}) by turning strong-form PDEs into weak forms,
$a(u,v)=\ell(v)$, drawn from a single four-operator family
$(Q,\diffop)$. We chose finite-element spaces for each differential operator and
saw them interlock as the de Rham complex (\cref{ch:functionspaces},
\cref{ch:derham}). We assembled the bilinear forms into sparse operators by a fold
over terms (\cref{ch:galerkin}, \cref{ch:assembly}), and we saw how that fold
runs matrix-free on the reference element (\cref{ch:refelement},
\cref{ch:matrixfree}). And in this chapter we performed the two finishing
operations---essential-BC elimination and divergence-free projection---that turn
the raw assembled system into a well-posed, physical one.

The result is a concrete linear-algebra problem: an operator $\mat{K}$, a
right-hand side $\vb$, perhaps a mass operator $\mat{M}$ for an eigenproblem, all
finished and ready. What we have \emph{not} done is solve it. Every analysis in
this book ends in a solve---a linear solve, an eigenvalue solve, a time march---and
those solves are where the simulator spends its time and where its deepest
structure lives. Part~III turns to the calculus that performs them: the
combinators and folds that express a Krylov iteration, a time integrator, an
eigenvalue sweep as compositions of pure, immutable-tensor operations. Part~II
built the problem; Part~III builds the machine that answers it.

\summarysection
Two finishing operations stand between the assembled system $\mat{K}\vx=\vb$ and
its solution, and both are \emph{separable post-compositions}---not steps inside
the assembly fold. \textbf{(1)~Essential boundary conditions.} The constrained
(Dirichlet) DoFs form an \emph{essential set} $E$; the rest are \emph{free}.
Imposing the condition partitions the system into free/essential blocks and
reduces it to the free block, $\mat{K}_{FF}\vx_F=\vb_F-\mat{K}_{FE}\vx_E$. In
place, \texttt{eliminate\_essential\_bc} zeroes the essential rows and columns and
sets the eliminated diagonal by a policy (\texttt{DIAG\_ONE} for the solve-side,
\texttt{DIAG\_ZERO} for the energy/mass-block), which is the free-block projection
$\mat{P}_F\mat{K}\mat{P}_F$ up to the chosen diagonal. A nonzero (inhomogeneous)
boundary value is handled by \emph{lifting}: \texttt{eliminate\_rhs} forms
$\vb\leftarrow\vb-\mat{K}\vx_{\mathrm{ess}}$, moving the known boundary
contribution into the right-hand side---a single matrix--vector product and an
\texttt{axpy}. Natural (Neumann) conditions, by contrast, enter during assembly as
boundary integrators and cost nothing when homogeneous. \textbf{(2)~Helmholtz
decomposition and divergence-free projection.} Any vector field splits uniquely as
$\vF=\vF_{\mathrm{df}}+\Grad\psi$ (divergence-free plus gradient), and at the
$\Hcurl$ slot of the de Rham complex the gradient part is exactly the kernel of
the curl. The discrete projector
$\mat{P}=\mat{I}-\mat{G}(\mat{G}^{\!\top}\mat{M}\mat{G})^{-1}\mat{G}^{\!\top}\mat{M}$
is an $\mat{M}$-orthogonal projection onto the divergence-free subspace; its apply
runs weak-divergence $\to$ essential-zeroing $\to$ inner Poisson solve $\to$
gradient subtract, and so is a constructed operator carrying an inner
\texttt{ksp\_solve}. The projection is needed because the curl--curl operator has
a huge gradient null space ($\mat{K}\Grad\phi=0$); without projection,
eigensolvers return spurious zero modes and iterative solvers stall. Both finishing
operations close the discretization that Part~II built.

\furtherreading
The variational treatment of boundary conditions---essential versus natural, the
trial/test-space constraint, and the lifting of inhomogeneous data---is developed
in any rigorous finite-element text; for the vector ($\Hcurl$, $\Hdiv$) case
specific to Maxwell problems, Monk~\citep{monk2003} is the standard reference, with
careful treatment of tangential-trace boundary conditions and the discrete spaces
on which they live. The Helmholtz decomposition and its de Rham reading are the
heart of finite-element exterior calculus; Arnold, Falk, and
Winther~\citep{arnold2006feec} give the definitive modern account, in which the
divergence-free projection is the Hodge decomposition at the relevant slot of the
discrete complex. The auxiliary-space and projection techniques that make the
curl--curl null space tractable for iterative solvers---the practical engine behind
the inner Poisson solve of \cref{sec:divfree}---are the subject of Hiptmair and
Xu~\citep{hiptmairxu2007}, whose AMS (auxiliary-space Maxwell solver) preconditioner
is built on exactly the discrete-gradient and mass operators that appear in
\cref{eq:projector}.

\exercisessection
\begin{enumerate}[nosep]
  \item \textbf{(Reading the partition.)} On the mesh of \cref{fig:partition}, the
  essential set $E$ is the boundary nodes and the free set $F$ is the interior.
  If the bottom edge is held at voltage $V$ and the top edge grounded ($0$), with
  the side edges insulating (homogeneous Neumann), which nodes are in $E$ and which
  in $F$? Which essential nodes carry $V$, which carry $0$, and what happens to the
  side-edge nodes?

  \item \textbf{(The two policies.)} Explain in one or two sentences why the
  \texttt{DIAG\_ONE} policy (essential diagonal $=\mat{I}_E$) is appropriate when
  the eliminated operator will be inverted by a linear solver, while
  \texttt{DIAG\_ZERO} (essential diagonal $=\mathbf{0}$) is used for the blocks of
  a generalized eigenvalue problem. What would go wrong if you used
  \texttt{DIAG\_ZERO} for the linear solve, or \texttt{DIAG\_ONE} for an eigenvalue
  block? (Hint: think about what eigenvalue an identity diagonal block contributes.)

  \item \textbf{(Lifting, by hand.)} Repeat \cref{wex:lift1d} but with both ends
  inhomogeneous: $u(0)=a$, $u(1)=b$. Form $\vx_{\mathrm{ess}}=(a,0,0,b)^{\!\top}$,
  compute $\mat{K}\vx_{\mathrm{ess}}$ and the lifted RHS
  $\vb'=\vb-\mat{K}\vx_{\mathrm{ess}}$, and solve the reduced free system. Confirm
  the answer is the linear interpolant $u(x)=a+(b-a)x$ sampled at the nodes.

  \item \textbf{(Separability.)} Suppose $\mat{K}=\mat{K}_1+\mat{K}_2$ is the sum
  of two assembled term-contributions (say a stiffness plus a mass term). Using the
  projection form \cref{eq:pfproj}, show that under the \texttt{DIAG\_ZERO} policy
  $\texttt{eliminate\_essential\_bc}(\mat{K}_1+\mat{K}_2,E)
  =\texttt{eliminate\_essential\_bc}(\mat{K}_1,E)+\texttt{eliminate\_essential\_bc}(\mat{K}_2,E)$.
  Then explain why the same distributivity fails for \texttt{DIAG\_ONE}, and state
  the corrected identity (the identity $\mat{I}_E$ is added once, not twice).

  \item \textbf{(Idempotence of the projector.)} Verify $\mat{P}^2=\mat{P}$ for the
  numerical instance of \cref{wex:projcheck} by direct matrix multiplication of
  $\mat{P}=\tfrac12\bigl[\begin{smallmatrix}1&-1\\-1&1\end{smallmatrix}\bigr]$. Then
  construct a second instance with $\mat{M}=\mathrm{diag}(2,1)$ (a nontrivial
  $\mat{M}$-inner-product) and the same
  $\mat{G}=\bigl[\begin{smallmatrix}1\\1\end{smallmatrix}\bigr]$: compute
  $\mat{G}^{\!\top}\mat{M}\mat{G}$, $\mat{S}$, and $\mat{P}$, and confirm $\mat{P}$
  is still idempotent and still annihilates $\mat{G}\psi$. (You will find $\mat{P}$
  is no longer symmetric---it is $\mat{M}$-orthogonal, not ordinary-orthogonal.)

  \item \textbf{(Why gradients are killed.)} Using only the identity
  $\mat{G}^{\!\top}\mat{M}\mat{G}=\mat{S}^{-1}$, give the two-line proof that
  $\mat{P}\,\mat{G}\psi=0$ for every $\psi$ (i.e.\
  $\mathrm{Ker}(\mat{P})\supseteq\mathrm{Range}(\mat{G})$). Then explain, in one
  sentence each, why this is exactly the statement that $\mat{P}$ removes the
  curl--curl null space, and why the inner solve in step~3 of
  \cref{fig:projapply} must be converged tightly for the cancellation to hold.

  \item \textbf{(Spurious modes.)} The curl--curl eigenproblem
  $\mat{K}\vx=\lambda\mat{M}\vx$ has every discrete gradient $\mat{G}\phi$ as an
  eigenvector. What is its eigenvalue $\lambda$? Explain why these modes cluster at
  the bottom of the spectrum and why an eigensolver seeking the lowest physical
  resonances is exactly the one most confused by them. How does projecting every
  iterate with $\mat{P}$ remove the problem?

  \item \textbf{(Natural versus essential, bookkeeping.)} A driven-cavity problem
  has a PEC wall (tangential $\vE=0$, essential), an impedance wall (a natural
  Robin-type condition contributing a boundary mass term), and a port (an
  inhomogeneous essential excitation). For each of the three, state: (i) whether it
  is imposed during assembly or as a post-composition, (ii) whether it touches the
  operator $\mat{K}$, the right-hand side $\vb$, or both, and (iii) whether a
  homogeneous version of it would cost anything. Organize your answer as a small
  table.
\end{enumerate}
